\subsection{Rational exactness of independently ordered binary vector spaces}

The rational entry for the ordered-vector example can be decided without
assuming a Ramsey theorem for two orders on vector spaces with arbitrary
linear embeddings.  We use a different category in which the embeddings
preserve a distinguished basis, and then transfer finite affine equalizers.

\begin{lemma}[Proved: Ramsey property with a distinguished ordered basis]
\label{A:ordered-basis-ramsey}
Let $\mathcal C$ have as objects finite positive-dimensional
$\mathbf F_2$-vector spaces $W$, a distinguished linearly ordered basis
$E_W$, and an independent linear order on $W\setminus\{0\}$.
Its arrows are injective linear maps taking basis elements to basis
elements and preserving both orders.  Then $\mathcal C$ is Ramsey,
has finite hom-sets and common successors, and is exact over every ring.
\end{lemma}
\begin{proof}
An object is equivalently a finite linearly ordered set $E$ with an
independent linear order on its nonempty subsets: a subset represents
the sum of its basis elements.  Fix objects $A,B$. If $A$ does not
embed into $B$, the Ramsey assertion is vacuous; otherwise
$|E_A|\leq |E_B|$. Put $b=|E_B|$.
For a nonempty subset $S=\{s_1<\cdots<s_k\}$ of size at most $b$, encode
$S$ by the $b$-tuple
\[
 \tau_b(S)=(s_1,\ldots,s_k,
                 \underbrace{s_k,\ldots,s_k}_{b-k\text{ entries}}).
\]
The distinct entries recover $S$, so this encoding is injective and
commutes with increasing injections.  On $E_A$ and $E_B$ put the
$b$-ary relation consisting of all such tuples.  Order that relation by
the given order on nonempty subsets, and retain the basis order as the
vertex order.

Hubi\v cka--Ne\v set\v ril's theorem on totally ordered structures
\cite[Theorem~4.46, arXiv v4]{DHubickaNesetril2019} gives a finite
Ramsey target $C_0$ for these two encodings.  Here ``totally ordered''
means that the vertices and the tuples of the relation carry independent
linear orders, both preserved by embeddings.  The quoted theorem applies
to one relation of arbitrary finite arity, including repeated entries.

In $C_0$, delete relation tuples that are not of the canonical form
$\tau_b(S)$, and add every missing canonical tuple.  Retain the order
on the remaining old relation tuples and extend it to the added tuples
arbitrarily.  This operation preserves every encoded copy of $B$:
on its vertex set all canonical tuples were already present and no
noncanonical tuple was present, because its map into $C_0$ was an
embedding.  Its two orders are unchanged.  Moreover every encoded
$A$-copy inside such a $B$-copy was already an $A$-copy in $C_0$.
Thus a coloring after completion restricts to the old $A$-copies, and
a monochromatic $B$ supplied by $C_0$ remains monochromatic afterward.

Finally give the nonempty subsets of the target basis of size greater
than $b$ arbitrary positions extending the order already defined on
the smaller subsets.  The resulting full powerset order gives an object
of $\mathcal C$ with the required Ramsey property.
There are finitely many basis injections between two objects.
Common successors are obtained by taking disjoint bases, merging each
of the two old orders, and placing the new mixed subset sums arbitrarily.
The Ramsey exactness theorem therefore gives exactness over every ring.
\end{proof}

\begin{lemma}[Proved: finite averaging through an exact category]
\label{A:exact-cover-averaging}
Let $u:\mathcal C\to\mathcal A$ be a functor, where $\mathcal C$ is
nonempty, has common successors, and is exact over a commutative ring $R$.
For $Y\in\mathcal A$ define
\[
 X_Y(C)=R[\mathcal A(Y,uC)].
\]
For a finite family $a_0,\ldots,a_k:P\to Q$, suppose there is
$\mu\in\operatorname{colim}_{\mathcal C}X_Q$ of augmentation one
whose images under every precomposition map
$a_i^*:\operatorname{colim}X_Q\to\operatorname{colim}X_P$ agree.
Then the family has an affine equalizer over $R$ in $\mathcal A$.
\end{lemma}
\begin{proof}
Let $D:X_Q\to X_P^k$ have components $a_i^*-a_0^*$.
Exactness of the colimit gives a surjection
\[
 \operatorname{colim}_{\mathcal C}\ker D
       \longrightarrow
 \ker\bigl(\operatorname{colim}_{\mathcal C}D\bigr).
\]
Lift $\mu$ to the left side.  A colimit element uses only finitely
many objectwise representatives.  Send their objects to one common
successor $C$ and add their images.  This yields
$t\in\ker D(C)\subseteq R[\mathcal A(Q,uC)]$ representing $\mu$.
Augmentation is natural, so $\operatorname{aug}t=1$.
The equality $D(C)t=0$ is precisely the desired affine equality.
\end{proof}

\begin{theorem}[Proved: rational exactness of the ordered-vector example]
\label{A:ordered-vector-rational-exactness}
Let $\mathcal A$ be the category of finite nonzero
$\mathbf F_2$-vector spaces with an arbitrary linear order on their
nonzero vectors and all order-preserving linear embeddings.
Then $\mathcal A$ is exact over $\Q$.
Consequently the example of
Theorem~\ref{A:ordered-vector-counterexample} is total, amenable, and
rationally exact, while it is not exact over $\mathbf F_2$ or over $\Z$.
\end{theorem}
\begin{proof}
Use the category $\mathcal C$ of
Lemma~\ref{A:ordered-basis-ramsey}, and let $u$ forget its distinguished
ordered basis.  First identify the components of the category of
elements of $C\mapsto\mathcal A(P,uC)$.
For an embedding $h:P\to uC$, let $U_h$ be the span of precisely those
distinguished basis vectors which occur with coefficient one in some
vector of $h(P)$.  Give $U_h$ the two induced orders.  Its isomorphism
type together with the marked map $h:P\to U_h$ is unchanged by arrows
of that category of elements.  Conversely $U_h\to C$ itself is such
an arrow.  Two marked closures of the same type therefore lie in the
same component, by their isomorphism and the two inclusions.  Thus
\[
 \operatorname{colim}_{\mathcal C}X_P
   =\Q[\{\text{isomorphism types of marked ordered basis closures of }P\}].
 \tag{BV.1}
\]
This uses only that the free-module functor preserves colimits.

Fix an arbitrary finite family $a_i:P\to Q$, and put
$p=\dim P$, $q=\dim Q$.  Let $Q^\vee$ be the binary linear dual and
form the binary simplex-code embedding
\[
 W=\mathbf F_2^{\,Q^\vee\setminus\{0\}},\qquad
 j:Q\longrightarrow W,\qquad
 j(v)=(\ell(v))_{\ell\in Q^\vee\setminus\{0\}}.
 \tag{BV.2}
\]
Its distinguished basis is the set of coordinate vectors, indexed by
the nonzero functionals $\ell$.  The map $j$ is injective because
linear functionals separate vectors.

Choose the basis order uniformly from all its possible orders.
Independently, order all vectors of
$W\setminus j(Q)$ uniformly, put that entire block before
$j(Q)\setminus\{0\}$, and put the latter block in the prescribed
order of $Q$.  Here and below the zero vector is omitted from all
orders.  Every such pair of choices gives an object $W_\omega$ of
$\mathcal C$ and an embedding $j:Q\to uW_\omega$.
There are finitely many choices; their uniform rational average defines
\[
 \mu=\frac{1}{|\Omega|}\sum_{\omega\in\Omega}[j:Q\to uW_\omega]
       \in\operatorname{colim}_{\mathcal C}X_Q,
 \qquad\operatorname{aug}\mu=1.
 \tag{BV.3}
\]

The basis closure $U_i$ of $ja_i(P)$ has coordinate set
\[
 E_i=\{\ell\in Q^\vee:\ell a_i\ne0\}.
\]
Every nonzero functional on $P$ occurs exactly $2^{q-p}$ times
among the restrictions $\ell a_i$, $\ell\in E_i$.  Consequently,
for any $i,j$ there is a coordinate bijection $E_i\to E_j$
intertwining the two marked embeddings of $P$.
We also have the exact intersection identity
\[
 U_i\cap j(Q)=ja_i(P).
 \tag{BV.4}
\]
Indeed, a vector $j(v)$ belongs to $U_i$ precisely when every
functional vanishing on $a_i(P)$ also vanishes on $v$; finite-dimensional
linear duality says exactly that $v\in a_i(P)$.

It follows from (BV.4) that the reference order induced on $U_i$
has all its vectors outside $ja_i(P)$ first, uniformly ordered,
followed by $ja_i(P)$ in the order of $P$.  The induced basis order
is independently uniform.  Restriction of a uniform order on a finite
set is uniform on each fixed subset; all fibres of the restriction
map have the same size.  The coordinate bijection just described
therefore shows that the marked ordered basis-closure type of $ja_i$
has exactly the same distribution for every $i$.  By (BV.1),
the images $a_i^*\mu$ all agree.

Lemma~\ref{A:exact-cover-averaging} now gives a rational affine
equalizer of the arbitrary family $a_i$.  Since $\mathcal A$ has
common successors, the all-family affine criterion proves rational
exactness.  The remaining properties in the statement were proved
for this same category in the earlier structural, amenability, and
parity results; no claim of integral exactness is made.
\end{proof}

\paragraph{Scope and source audit.}
The Ramsey input is Theorem~4.46 in the opened arXiv v4 of
\cite{DHubickaNesetril2019}, dated 3 September 2019:
the definitions are on printed page 75 and the theorem and complete
incidence-closure proof on printed page 76.  The padded-tuple
completion above verifies the passage to powerset orders explicitly.
The finite averaging proof uses an auxiliary exact category and
finite supports; it does not assert that its basis expansion is
precompact or invoke a general theorem about free order expansions
of vector spaces.  The universal implications in A.3(iii) and the
total-category version of C.4 remain undecided by this result.
